%% ma: L11-B09-0019
%% lop: 11
%% chuong: 3
%% bai: 9
%% dang: 11.9.4
%% loai: TL
%% muc: TH
%% dap_an: "-"
%% nguon: "de-on-tap-giua-ki-1-toan-11-nam-2026-2027-truong-thpt-ngo-quyen-da-nang.pdf (M09, Đề ôn tập giữa kì 1 Toán 11 THPT Ngô Quyền, Đà Nẵng), Mã đề 0114, Phần IV Câu 3"
%% kien_thuc: "Bài 9 (các số đặc trưng đo xu thế trung tâm)"
%% trang_thai: chinh-thuc
\begin{ex}
Trong một nghiên cứu, các nhà khoa học ghi lại tuổi thọ của một loài côn trùng và cho kết quả bởi bảng sau:
\begin{center}\begin{tabular}{|l|c|c|c|c|c|}\hline
Tuổi (ngày) & $[0;20)$ & $[20;40)$ & $[40;60)$ & $[60;80)$ & $[80;100]$\\\hline
Số lượng & $6$ & $12$ & $23$ & $31$ & $39$\\\hline
\end{tabular}\end{center}
Hãy cho biết tuổi thọ trung bình của loài côn trùng trên (kết quả làm tròn đến hàng phần mười).
\loigiai{$n=111$, giá trị đại diện $10;30;50;70;90$. \[\bar x=\dfrac{10\cdot6+30\cdot12+50\cdot23+70\cdot31+90\cdot39}{111}=\dfrac{7250}{111}\approx65{,}3\text{ (ngày).}\]}
\end{ex}
